\documentclass[11pt,reqno]{amsart}
\usepackage[utf8]{inputenc}
\usepackage{amsmath,amssymb,amsthm,mathtools}
\usepackage[margin=1in]{geometry}
\usepackage[colorlinks=true,linkcolor=blue,citecolor=blue,urlcolor=blue]{hyperref}

\theoremstyle{plain}
\newtheorem{theorem}{Theorem}
\newtheorem{proposition}[theorem]{Proposition}
\newtheorem{lemma}[theorem]{Lemma}
\newtheorem{corollary}[theorem]{Corollary}
\theoremstyle{definition}
\newtheorem{definition}[theorem]{Definition}
\newtheorem{remark}[theorem]{Remark}

\theoremstyle{plain}
\newtheorem*{structthm}{Structure theorem (unsigned kernel; certified perimeter)}
\newtheorem*{lemA}{Lemma A (modes are cycles)}
\newtheorem*{lemB}{Lemma B (cone split at the least vertex)}
\newtheorem*{lemC}{Lemma C (graded reduction)}
\newtheorem*{RDL}{Realizable Decomposition Lemma (the crux)}
\newtheorem*{onestep}{One-step elimination lemma}
\theoremstyle{definition}
\newtheorem*{hypLG}{Hypothesis (LG)$_{v,a}$}

\newcommand{\rank}{\operatorname{rank}}
\newcommand{\elem}{\operatorname{elem}}
\newcommand{\ru}{\operatorname{ru}}
\newcommand{\cpc}{\operatorname{cap}}
\newcommand{\Z}{\mathbb{Z}}
\newcommand{\Q}{\mathbb{Q}}
\newcommand{\F}{\mathbb{F}}
\newcommand{\eO}{\mathcal{O}}
\newcommand{\Zbal}{Z_{\mathrm{bal}}}
\newcommand{\dbu}{\partial^{u}}
\newcommand{\dd}{\delta}
\newcommand{\st}{\textsc}

\begin{document}

\title[A balance-mode structure theorem for the unsigned kernel]
{A balance-mode structure theorem for the unsigned kernel\\of the sphenic window complex}
\author{Lee Rich}
\thanks{ORCID: \href{https://orcid.org/0009-0004-8559-5177}{0009-0004-8559-5177}.
The Improbability Lab. \emph{Working manuscript, 2 July 2026.}}
\email{leemjrich@gmail.com}
\date{2 July 2026}

\begin{abstract}
Let $C$ be the sphenic window complex at level $N$: the $2$-complex whose triangles are the
squarefree three-prime integers $pqr$ in $W=(N/2,N]$, with edges the co-occurring prime pairs
and vertices the primes. Let $\dbu$ be its \emph{unsigned} boundary, the all-positive incidence
map, which presents $C$ as an oriented hypergraph with every adjacency negative in the sense of
Rusnak. We describe the rational kernel of $\dbu$ as a span of explicit \emph{balance modes}
fibered over vertex-pairs, and certify it at an explicit finite perimeter. For the five bundled
complexes $N\in\{50000,100000,200000,300000,500000\}$, certified over $\Q$ and an eight-prime
battery at all five $N$ (and over every prime characteristic outside $\{2,3\}$ on the two-stage
integer-certified subset $N\in\{10^5,2\times10^5\}$),
\[
\ker\dbu=\operatorname{span}\bigl\{\,m_{a,b,z}=\textstyle\sum_g z_g([a\cup g]-[b\cup g]):
a<b\text{ prime},\ z\in\Zbal(L_{a,b})\,\bigr\},
\]
where $L_{a,b}$ is the common link graph of the pair and $\Zbal$ its balance kernel (the
Zaslavsky even-circle cycle space). Consequently, at that perimeter, $e:=\dim\ker\dbu$ is a
fibered balance count, $\ru=F-e$, the frustration term is $\dd=\elem-e$, and the canonical
density $\kappa=(b_1-\dd)/V$ is determined by link-graph balance and connectivity combinatorics.
The characteristics $2$ and $3$ are exceptional and exactly understood: over $\F_2$ the kernel
is $\elem$, over $\F_3$ it is $e+1$. The generators are the classified balanced and balanceable
circuits of the oriented-hypergraph program of Rusnak \cite{Rusnak2013} and of
Rusnak--Li--Xu--Yan--Zhu \cite{RLXYZ2022}; the apparently unstated \emph{spanning} statement
addresses their declared open problem. We prove the reduction unconditionally, reduce the
remaining content to a single finite leakage-gauge condition (LG), and certify (LG) at the
perimeter above. The all-$N$ form is \st{[open]}, with the band localization proved and a
richness program stated.
\end{abstract}

\maketitle

\section{The unsigned kernel, and the statement at the certified perimeter}\label{sec:stmt}

\subsection{The object}
The complex $C$ is as in the companion Paper~3: triangles are the sphenics $pqr\in(N/2,N]$, edges
the co-occurring prime pairs, vertices the primes. The signed boundary $\partial_2$ has rank
$F-\elem$ by the slab identity. This paper is about the \emph{unsigned} boundary
\[
\dbu[p,q,r]=[q,r]+[p,r]+[p,q]\qquad(\text{all coefficients }+1),
\]
the incidence map of $C$ as an oriented hypergraph with all-positive incidence, equivalently all
adjacencies negative in Rusnak's dictionary. Over a field of characteristic $\ne2$ the two
boundaries differ, and their nullities differ by the frustration term
\[
\dd:=\ru-\rank\partial_2=(F-e)-(F-\elem)=\elem-e,\qquad e:=\dim\ker\dbu,\ \ \ru:=\rank\dbu.
\]
The term $\dd$ is the one quantity in the whole coprime-graph arc that has resisted a rank-free
closed form; the point of the structure theorem is to express $e$, hence $\dd$, as a balance
count.

\subsection{The fiber objects}
For a vertex $v$, the \emph{link} $L_v$ is the graph on primes with an edge $\{x,y\}$ whenever
$\{v,x,y\}$ is a triangle of $C$. For a prime pair $\{a,b\}$, the \emph{common link} $L_{a,b}$
is the graph whose edges $g=\{x,y\}$ satisfy that both $\{a,x,y\}$ and $\{b,x,y\}$ are triangles.
Signing every adjacency negative (equivalently, the all-positive incidence), the \emph{balance
kernel} $\Zbal(G)$ is the Zaslavsky even-circle cycle space \cite{Zaslavsky}: the space spanned
by the even circles and the odd-circle barbells of $G$, the cycle space of the signed graphic
matroid of the all-negative signing. For each pair and each $z\in\Zbal(L_{a,b})$, the
\emph{balance mode} is $m_{a,b,z}=\sum_g z_g([a\cup g]-[b\cup g])$, with $g$ ranging over the
edges of $L_{a,b}$.

\subsection{The theorem, scoped}
\begin{structthm}
\textnormal{\st{[banked as C-0746 at the certified perimeter]}}\ \ For each of the five bundled
complexes $N\in\{50000,100000,200000,300000,500000\}$, certified over $\Q$ and an eight-prime
battery at all five $N$, and over every prime characteristic $p\notin\{2,3\}$ on the two-stage
integer-certified subset $N\in\{10^5,2\times10^5\}$,
\[
\ker\dbu=\operatorname{span}\{\,m_{a,b,z}: a<b\text{ prime},\ z\in\Zbal(L_{a,b})\,\}.
\]
\end{structthm}
\noindent We state the scope plainly. The inclusion $\supseteq$ (the modes are kernel vectors)
is an unconditional lemma for all $N$ (\S\ref{sec:core}); the reduction of $\subseteq$ to a
finite condition (LG) is unconditional for all $N$ (\S\ref{sec:peel}). What is verified rather
than proved in general is (LG) itself, verified at the five $N$ by an exact certificate stack
(\S\ref{sec:cert}). The characteristic scope matches the certificate inventory exactly: the
leakage-gauge two-rank test over an eight-prime battery covers all five $N$ and pins the kernel
over $\Q$, while the all-characteristic closure (no bad-prime window) is certified by the
two-stage integer certificates only on $N\in\{10^5,2\times10^5\}$; we do not claim all
characteristics at the three larger $N$, where those integer certificates are not yet in the
repository. The all-$N$ statement is open (\S\ref{sec:allN}), and every downstream consequence
inherits the certified-perimeter scope.

\section{The proved core: modes are cycles, and the cone split}\label{sec:core}

\begin{lemA}
\textnormal{\st{[proved, all $N$]}}\ \ For any prime pair $\{a,b\}$ and any
$z\in\Zbal(L_{a,b})$ (zero unsigned vertex sums on $L_{a,b}$), $\dbu m_{a,b,z}=0$.
\end{lemA}
\begin{proof}
Examine each facet coefficient of $\dbu m$. A facet $\{x,y\}$ interior to an edge $g$ receives
$z_g$ from $[a\cup g]$ and $z_g$ from $[b\cup g]$ with the mode's sign difference, so the
coefficients cancel. A facet $\{a,u\}$ receives $\sum_{g\ni u}z_g$, the unsigned vertex sum of
$z$ at $u$, which is zero by balance; likewise $\{b,u\}$. The facet $\{a,b\}$ never occurs.
Hence $\dbu m=0$.
\end{proof}

Filter the kernel by least vertex. For a vertex $v$ let $C_{\ge v}$ be the subcomplex on
vertices $\ge v$, and let $w\in\ker\dbu|_{C_{\ge v}}$ have least vertex $v$. Split
$\operatorname{supp}(w)=A\sqcup B$, $A$ the cells containing $v$, and write $w_A$ for the tail
chain of the $A$-cells.
\begin{lemB}
\textnormal{\st{[proved, all $N$; characteristic $\ne2$]}}\ \ \emph{(i)}
$w_A\in\Zbal(L_v^{\ge v})$ (the $v$-facet block); \emph{(ii)} $w_A=-\dbu(u)$ exactly, where $u$
is the $B$-part of $w$ (the $v$-free block, both sides supported on $L_v$-edges); \emph{(iii)}
if $B=\varnothing$ then $w=0$.
\end{lemB}
\noindent Writing $W_v$ for the space of $A$-parts of kernel vectors of $C_{\ge v}$ with least
vertex $v$, we obtain $W_v=\Zbal(L_v^{\ge v})\cap\dbu(\text{chains of }C_{>v})$. The second
condition is not decorative: the naive version without it is false, with a measured deficit of
$1276$ at $N=10^5$; any proof ignoring the boundary witness is refuted by that counterexample.

\begin{lemC}
\textnormal{\st{[proved, all $N$]}}\ \ The $A$-part map identifies the graded quotient
$K(C_{\ge v})/K(C_{>v})$ with $W_v$, and the modes with $a=v$ have $A$-parts exactly
$\sum_{b>v}\Zbal(L_{v,b}^{\ge v})\subseteq W_v$. Hence, by downward induction on the finitely
many least-vertex levels (primes up to roughly $N^{1/3}$), the structure theorem for $C_{\ge v}$
follows from the theorem for $C_{>v}$ together with the following.
\end{lemC}
\begin{RDL}
$W_v\subseteq\sum_{b>v}\Zbal(L_{v,b}^{\ge v})$ at every level: every balanced, link-supported
boundary of the upper subcomplex splits into pair-fiber balance cycles.
\end{RDL}
\noindent Lemmas A, B and C are unconditional; the whole remaining content is the crux, which
\S\ref{sec:peel} reduces to a finite gauge condition.

\section{The peeling architecture, with corrected bookkeeping}\label{sec:peel}

We prove the crux under a finite hypothesis (LG). The essential correction over earlier attempts
is the bookkeeping: one does \emph{not} hold the witness boundary fixed while removing a layer;
each peeling step subtracts a certified element $q$ from the current boundary, so the boundary
itself changes as $h\mapsto h-q$.

\subsection{Notation}
Order vertices by prime value. For $a>v$ let $X_a$ be the cells with all vertices $\ge a$, and
$a^+$ the next vertex. For $b\ge a$ let $L_{v,b}^a$ be the common link on vertices $\ge a$ with
$b$ removed, and $T_v^a=\sum_{b\ge a}\Zbal(L_{v,b}^a)$, so $T_v^{a_0}=T$, the target pair-fiber
span ($a_0$ the first vertex above $v$). Let $L_a$ be the link of $a$ inside $X_a$ and
$T_a=\sum_{c>a}\Zbal(L_{a,c})$. For $u\in C_2(X_a)$, its $a$-tail $\ell_a(u)$ keeps the tail
$[x,y]$ of each cell $[a,x,y]$ with least vertex $a$ and sends all other cells to $0$.

\subsection{Mode halves and the leakage tail}
For $z\in\Zbal(L_{v,b}^a)$, the $C'$-half of the full $v$-mode is
$R_{v,b}(z)=-\sum_g z_g[b\cup g]$, and balance gives $\dbu R_{v,b}(z)=-z$ on the link edges.
Summing over $b$, for $\mathbf z=(z_b)$ one gets $R(\mathbf z)$ with $\dbu R(\mathbf z)=-q(\mathbf
z)$, where $q(\mathbf z)=\sum_b z_b\in T_v^a$. The \emph{leakage-tail map} is
$\beta_{v,a}(\mathbf z)=\ell_a(R(\mathbf z))$, an explicit linear map into $C_1(L_a)$: the tail
edge $\{x,y\}$ receives $-(z_a)_{xy}-(z_x)_{ay}-(z_y)_{ax}$, absent terms read as $0$.

\subsection{Straddle localization by the window (P1)}
Let $u\in C_2(X_a)$ with $h=\dbu u\in\Zbal(L_v^a)$ and $t=\ell_a(u)$. The imbalance of $t$ at
$x$ equals $h_{ax}$ when $\{a,x\}$ is a link edge and $0$ otherwise, so it is supported on the
straddle set $S_{v,a}=\{x>a:\{a,x\}\in L_v\}$ and sums to zero there. The window supplies (P1):
if a cell $\{a,x,y\}$ shares the edge $\{a,x\}$ with a cell $\{v,a,x\}$, their third vertices
$y,v$ differ by a factor below $2$, so $y<2v$. In particular, once $a\ge2v$ no straddling
contribution occurs and the $a$-tail is automatically balanced, confining the difficulty to the
factor-$2$ band $a\in(v,2v)$.

\subsection{The finite leakage-gauge hypothesis}
Let $\mathcal U_{v,a}=\{\ell_a(u): u\in C_2(X_a),\ \dbu u\text{ supported on and balanced in
}L_v^a\}$.
\begin{hypLG}
For every $t\in\mathcal U_{v,a}$ there is $\mathbf z$ with $t+\beta_{v,a}(\mathbf z)\in T_a$;
equivalently, in $C_1(L_a)/T_a$, $\overline{\mathcal U}_{v,a}\subseteq\operatorname{im}
\overline{\beta}_{v,a}$. A stronger link-only sufficient form is
$C_1(L_a)=\operatorname{im}\beta_{v,a}+T_a$.
\end{hypLG}

\subsection{One-step elimination and the induction}
\begin{onestep}
\textnormal{\st{[proved under (LG)$_{v,a}$]}}\ \ Given $u\in C_2(X_a)$ with
$h=\dbu u\in\Zbal(L_v^a)$, there exist $q\in T_v^a$ and $u^+\in C_2(X_{a^+})$ with
$\dbu u^+=h-q$ and $h-q\in\Zbal(L_v^{a^+})$.
\end{onestep}
\begin{proof}
Let $t=\ell_a(u)$. By (LG)$_{v,a}$ pick $\mathbf z$ with $p:=t+\beta_{v,a}(\mathbf z)\in T_a$,
and set $R=R(\mathbf z)$, $q=q(\mathbf z)\in T_v^a$, so $\dbu R=-q$ and
$\ell_a(R)=\beta_{v,a}(\mathbf z)$. Since internal kernel tails are exactly $T_a$ (Lemma~C inside
$X_a$), choose an internal cycle $N$ with $\ell_a(N)=p$. Put $u^+=u+R-N$. Then $\dbu u^+=h-q$,
and $\ell_a(u^+)=t+\beta_{v,a}(\mathbf z)-p=0$, so $u^+$ has no cell with least vertex $a$ and
lies in $C_2(X_{a^+})$. Finally $h-q$ is balanced and, being $\dbu u^+$ with $u^+$ free of $a$,
has no edge incident to $a$, so $h-q\in\Zbal(L_v^{a^+})$.
\end{proof}
Reverse induction on $a$ proves: if $u\in C_2(X_a)$ and $\dbu u=h\in\Zbal(L_v^a)$, then
$h\in T_v^a$. The base after the largest vertex is trivial; the step is the one-step lemma plus
the induction hypothesis on $h-q$. Taking $a=a_0$ gives $W_v\subseteq T$, the crux, hence the
structure theorem under the (LG) hypotheses. Unwinding exhibits the corrected bookkeeping: chains
$u_0=u,u_1,\dots,u_r=0$ and certified $q_i\in T$ with $\dbu u_i=h_i$, $h_{i+1}=h_i-q_i$, so
$h=q_0+\dots+q_{r-1}\in T$. The witness boundary is not held fixed; each peel subtracts a
certified pair-fiber element.

\section{(LG) as a finite condition, and the certificate system}\label{sec:cert}

\subsection{The finite rank form}
For fixed $a$, (LG)$_{v,a}$ is a single rank equality. Let $U_a$ span the occurring tails
$\mathcal U_{v,a}$ (a kernel of an explicit two-block boundary matrix), $G_a$ span
$\operatorname{im}\beta_{v,a}$, and $H_a$ span $T_a$. Then
\[
\text{(LG)}_{v,a}\text{ holds}\iff\rank[\,G_a\ H_a\ U_a\,]=\rank[\,G_a\ H_a\,],
\]
a machine-checkable equality of ranks of explicit integer matrices, one per level pair. The band
localization confines the nontrivial checks to $a\in(v,2v)$.

\subsection{The certificate stack}
The structure theorem is certified at the five bundled complexes by the following, all
reproducible from the repository root (bundles under \texttt{certificates/} are gzipped JSON with
SHA-256 manifests).
\begin{itemize}
\item \emph{Leakage-gauge rank certificates.} The two-rank test passes at $616$ level pairs
across the five $N$, integer-exact. \st{[certified]}
\item \emph{Eight-prime battery at emission.} Every rank confirmed over eight prime
characteristics as certificates are emitted. \st{[certified]}
\item \emph{Modular sweep.} A $32{,}868$-check sweep over characteristics up to $997$ finds no
bad prime below $1000$. \st{[verified]}
\item \emph{Two-stage integer certificates.} For $N\in\{10^5,2\times10^5\}$ ($198$ pairs), a
unit-pivot unimodular reduction followed by an exact Smith normal form of the residual block
produces integer invariant factors with none divisible by any prime outside $\{2,3\}$, closing
the field story on that subset without a bad-prime window above $1000$; at the three larger $N$
the certificates cover $\Q$ and the eight-prime battery, not all characteristics. \st{[certified
on $N\in\{10^5,2\times10^5\}$]}
\item \emph{Independent enumeration replay.} A separate program rebuilds every cell inventory
from $N$ alone and reproduces the fiber and generator data. \st{[verified clean]}
\item \emph{Independent checker replay.} A second implementation, using spanning-tree balance
bases and a different elimination order, re-verifies every rank over the battery. \st{[verified
clean]}
\end{itemize}
The formal adjudication banked the $\Q$ plus eight-prime perimeter at $N\in\{10^5,2\times10^5\}$
($198$ pairs, $1584$ rank-pair checks); the five-$N$, $616$-certificate extension is on the same
epistemic footing, a theorem modulo a hashed finite certificate.

\subsection{The bug that failed loudly (reproducibility)}
The independent checker earned its keep. Its first version had a sign error in the barbell
balance-closure, and it failed loudly: the checker disagreed with the emitter on $22$ of $78$
fibers, at all eight primes at once. The fault was in the checker, not the certificates, and a
loud, characteristic-independent disagreement is exactly what an independent replay should
produce when something is wrong. The checker was rewritten with a brute-forced sign closure plus
a per-component dimension guard (each balance component contributes $E-V+[\text{bipartite}]$),
after which emitter, checker and enumerator agree on every pair. A certificate stack whose
independent replays never disagree has not really been tested; the loud disagreement is the
evidence that the replay is genuinely independent, and the adjudication credited it as such.

\section{Characteristic stratification}\label{sec:strat}

The kernel dimension of $\dbu$ is not field-independent, and the two exceptional characteristics
are exactly the ones the oriented-hypergraph taxonomy predicts.
\begin{itemize}
\item \emph{Characteristic $2$} \st{[proved; from C-0744]}: over $\F_2$ the unsigned and signed
boundaries coincide, so $\ker\dbu=\ker\partial_2$ has dimension $\elem$ (the fitting-cell space,
not the balance-mode space).
\item \emph{Characteristic $3$} \st{[verified; certified on the perimeter]}: over $\F_3$ the
kernel dimension is exactly $e+1$, one extra cross-theta class. This is the all-entrant
$\mathrm{GF}(k)$ phenomenon: an all-entrant minimal $k$-cross-theta is minimally dependent only
over $\mathrm{GF}(k)$, so at $k=3$ exactly one such dependency becomes rational over $\F_3$ and
disappears elsewhere. The measured $e+1$ is a confirmed prediction of Rusnak's characteristic
theory.
\item \emph{Every other characteristic} \st{[certified on the perimeter]}: over $\Q$ and over
every prime outside $\{2,3\}$ the kernel dimension is $e$ (two-stage subset; battery-scoped
elsewhere), by the certificate stack of \S\ref{sec:cert}.
\end{itemize}
On the certified set the picture is a clean stratification: $\F_2$ gives $\elem$, $\F_3$ gives
$e+1$, and the generic $e$ holds everywhere else, with $\dd=\elem-e$ the characteristic-$2$ jump.

\section{Positioning in the oriented-hypergraph program}\label{sec:pos}

The generators are classified circuits; only the spanning statement is ours to claim, and even
that is scoped. Under the standard dictionary (all-positive incidence corresponds to
all-negative adjacency, so a positive circle is an even circle and Zaslavsky balance is
evenness), the two mode families are known objects. The parity boxes are \emph{balanced
circuits} in the sense of Rusnak, \emph{Oriented Hypergraphs: Introduction and Balance}
\cite{Rusnak2013} (the arXiv preprint \texttt{1210.0943} carries the title \emph{Oriented
Hypergraphs I}); their classification of balanced circuits is the relevant result (Theorem 7.2.7
in the arXiv numbering, to be confirmed against the published version before quoting the number).
The barbell modes are \emph{balanceable circuits} in the sense of Rusnak, Li, Xu, Yan and Zhu,
\emph{Oriented Hypergraphs: Balanceability} \cite{RLXYZ2022}. We cite Lucas J.\ Rusnak as sole
author of the 2013 paper.

What is apparently unstated is the spanning statement. No nullity or spanning theorem exists in
the oriented-hypergraph program: the unbalanceable circuit family is their declared open problem,
and no global kernel-dimension formula for hypergraph incidence matrices is known at all. Our
theorem, restricted to the sphenic/window family and at the certified perimeter, asserts that the
two classified families already span the rational kernel, fibered over vertex-pairs by link
balance; equivalently, the unbalanceable circuits are rationally redundant on this family. There
is independent structural support: the all-entrant minimal $k$-cross-thetas are minimally
dependent only over $\mathrm{GF}(k)$, so the simplest unbalanceable circuits contribute no
rational kernel vectors, exactly the mechanism visible as the $\F_3$ stratum of \S\ref{sec:strat}.
We frame the contribution as an apparently unstated spanning equality over their taxonomy, not as
a new circuit theory. The fiber object is the Zaslavsky even-circle matroid, suggesting a
matroid-union reading of $e$ (a route, not a result). Adjacent but distinct is the Nagel--Reiner
complex-of-boxes machinery \cite{NagelReiner}, a different functor on the same skew objects; the
one named residual specialist check is Duval's relative Laplacian spectral recursion
\cite{DuvalLap}.

\section{The all-$N$ problem, and what is proved toward it}\label{sec:allN}

The all-$N$ structure theorem is \st{[open]}. The obstruction is precisely a uniform (LG). Two
things are proved toward it. \emph{Band localization} \st{[proved]}: by (P1), all straddle
imbalance lives in the factor-$2$ band, so for a level $a\ge2v$ every occurring $a$-tail is
automatically balanced and the nontrivial gauge condition is confined to partners $a\in(v,2v)$;
above $2v$ the peeling is straddle-free and unconditional. \emph{The richness program}
\st{[open, sharply stated]}: what remains is to show the leakage map is surjective onto the
occurring band tails uniformly in $N$, plausibly under an explicit finite richness hypothesis on
the band fibers (enough independent even circles and barbells in each $L_{v,b}$ with $v<b<2v$ to
realize every occurring straddle correction). This is a concrete, machine-testable family of
conditions; proving it uniformly, or replacing it by a band-local witness-modification argument,
would upgrade the theorem from the certified perimeter to all $N$.

\section{Consequences at the certified perimeter}\label{sec:conseq}

At the certified perimeter the whole coprime-graph arc becomes counting. The unsigned nullity
$e$ is the dimension of a sum of pair-fiber balance spaces, computed by inclusion-exclusion over
overlapping pairs; each component's balance dimension is $E-V+[\text{bipartite}]$, with the
non-bipartite fraction growing with scale (measured by the bipartite census at three of the
certified $N$: $0.73,0.81,0.88$), confining corrections to the completer fringe. The frustration
term $\dd=\elem-e$ becomes the characteristic-$2$ jump of a certified quantity, and the canonical
density $\kappa=(b_1-\dd)/V$ is reduced, at these $N$, to link-graph balance and connectivity
combinatorics. Combined with the exact signed profile of the companion Paper~3, this removes
the last uncertified rank from the sphenic ($\omega=3$) layer of the author's
$-1$-multiplicity decomposition \cite{Rich2026a,Rich2026b} at the perimeter; the present paper
is the fourth in that series.
The transcendence question for $\kappa$ at infinity remains tied to the all-$N$ richness program.

\appendix
\section{Machine certificates (reproducible)}
All scripts run from the repository root (concept DOI \texttt{10.5281/zenodo.21135221}).
\texttt{mode\_completeness.py} (and a $5\times10^5$ variant): global completeness, rank of the
mode family equals $e$ at $N\in\{5\times10^4,10^5,2\times10^5,5\times10^5\}$, values
$482,1270,3082,9402$. \texttt{graded\_completeness.py}: the graded crux verified level by level.
\texttt{atomic\_space\_test.py}: the atomic realizable decomposition at space level.
\texttt{lg\_rank\_test.py}, \texttt{twostage\_certify.py}: the (LG) two-rank and the two-stage
integer certificates. \texttt{lg\_replay\_check.py}, \texttt{enum\_replay.py}: the independent
checker and enumeration replays. \texttt{gauge\_anatomy2.py}, \texttt{slab\_gf2\_sweep.py}: the
$\F_3$ and $\F_2$ characteristic strata.

\section*{Acknowledgement / disclosure}
The balance-mode reformulation of the unsigned kernel and the reduction of completeness to the
finite leakage-gauge condition (LG) are the author's; the one-step elimination architecture of
\S\ref{sec:peel} originated in an AI-assisted solver round and was independently cold-refereed.
The manuscript was prepared with assistance from a large language model (Claude, by Anthropic)
for organisation, the classical background and prior-art positioning, the balance formalism, and
extensive symbolic and numerical verification; all statements, proofs, and their interpretation
are the author's and were checked independently, with banking-grade claims filed through the
laboratory's adversarial adjudication protocol before being marked proved. Results carry status
labels: \st{[proved]}, \st{[banked]}, \st{[verified]}, \st{[certified]}, \st{[open]}.

\begin{thebibliography}{9}
\bibitem{Rich2026a} L.~Rich, \emph{The $-1$-eigenspace of the coprime graph of integers:
prime-cube modes and an exact multiplicity}, Zenodo (2026), doi:10.5281/zenodo.20933665.
\bibitem{Rich2026b} L.~Rich, \emph{The $-1$-eigenvalue multiplicity of the coprime graph:
a conditional sublinear law}, Zenodo (2026), doi:10.5281/zenodo.21169867.
\bibitem{Rusnak2013} L.~J.~Rusnak, \emph{Oriented Hypergraphs: Introduction and Balance},
Electron.\ J.\ Combin.\ \textbf{20}(3) (2013), \#P48; arXiv:1210.0943.
\bibitem{RLXYZ2022} L.~J.~Rusnak, S.~Li, B.~Xu, E.~Yan, S.~Zhu, \emph{Oriented Hypergraphs:
Balanceability}, Discrete Math.\ \textbf{345} (2022); arXiv:2005.07722.
\bibitem{Zaslavsky} T.~Zaslavsky, \emph{Signed graphs}, Discrete Appl.\ Math.\ \textbf{4} (1982),
no.~1, 47--74.
\bibitem{DuvalLap} A.~M.~Duval, \emph{A relative Laplacian spectral recursion}, Electron.\ J.\
Combin.\ \textbf{11}(2) (2006), \#R26, math/0507130.
\bibitem{NagelReiner} U.~Nagel, V.~Reiner, \emph{Betti numbers of monomial ideals and shifted
skew shapes}, Electron.\ J.\ Combin.\ \textbf{16}(2) (2009), \#R3, arXiv:0712.2537.
\bibitem{Banerjee2025} S.~Banerjee, \emph{On structural properties and adjacency spectrum of
coprime graph of integers}, Asian-Eur.\ J.\ Math.\ (2025), doi:10.1142/S179355712550069X;
arXiv:2506.10583.
\end{thebibliography}

\end{document}
